%% kapa: exam=1 round=25 year=2014 period=1 form=A subject=civil_law subject_ko=민법 seq=2 number=2
%% source: Q-Net 감정평가사 기출문제 (cst00309) artlSeq=5101207 file=extracted/1008/5101207/2014년 제25회 감정평가사 1차시험 기출문제/1교시(A형).hwp
%% conversion: kordoc 4.24.51 -> tools/gen1.py (auto); review=proofread-sonnet; flags=none
\begin{question}{2}
신의성실의 원칙에 관한 설명으로 옳은 것은? (다툼이 있으면 판례에 의함)
\begin{choices}
  \item 강행법규 위반임을 알면서 법률행위를 한 자가 후에 그 법률행위가 강행법규 위반으로 인하여 무효임을 주장하는 것은, 특별한 사정이 없는 한, 신의칙에 반한다.
  \item 본인의 지위를 단독으로 상속한 무권대리인이 본인의 지위에서 상속 전에 행한 무권대리행위의 추인을 거절하여도 신의칙에 반하지 않는다.
  \item 사정변경 원칙에서의 ‘사정’은 계약의 기초가 된 일방 당사자의 주관적 사정을 말한다.
  \item 토지의 소유권을 적법하게 취득한 새로운 권리자에게 실효의 원칙을 적용함에 있어서는 종전 토지소유자가 자신의 권리를 행사하지 않았다는 사정을 고려하여야 한다.
  \item 법원은 당사자의 주장이 없더라도 직권으로 신의칙에 반하는 것임을 판단할 수 있다.
\end{choices}
\end{question}
\answer{5}
